\begin{center}
  \textbf{MỞ ĐẦU}
\end{center}
% \changefontsizes[16pt]{12pt}
\textbf{Lý do chọn đề tài}

Hội chứng ngưng thở khi ngủ do tắc nghẽn (Obstructive Sleep Apnea - OSA) được
đặc trưng bởi sự tắc nghẽn một phần hoặc hoàn toàn của đường hô hấp trên. Các
sự kiện này được xác định kéo dài ít nhất 10 giây, trong khi các nỗ lực hô hấp
vẫn tiếp tục thông qua cử động của lồng ngực và bụng.
% \cite{Epstein2009}
% Trong nghiên cứu \cite{Benjafield2019}, tác giả chỉ ra 
Hiện nay, có gần 1 tỷ người mắc OSA trên thế giới và theo GS.BS. Dương Quý Sỹ,
Việt Nam có tới 8.5\% người trưởng thành mắc OSA.
% \cite{nguoimacOSA_VN}. 
% OSA có ảnh hưởng lớn tới sức khỏe cả thể chất lẫn tinh thần đối với người mắc.
% Xia Wang trong nghiên cứu của mình có tới 25.760 người tham gia đã kết luận
% rằng: 
Có một số nghiên cứu chỉ ra rằng mỗi khi chỉ số ngưng - giảm thở tăng thêm 10
đơn vị, nguy cơ mắc bệnh tim mạch sẽ tăng 17\%.
% \cite{Wang2013_tim}. 
% Như trong nghiên cứu\cite{Salari2025}, nhóm tác giả đã
% phân tích dữ liệu từ 15 nghiên cứu với tổng cộng 42.924 người và kết quả cho
% thấy 
Ngoài ra, tỷ lệ mắc chứng buồn ngủ ban ngày quá mức (Excessive Daytime
Sleepiness - EDS) ở bệnh nhân ngưng thở tắc nghẽn khi ngủ trên toàn cầu là
khoảng 39.9\%.

Hội chứng ngưng thở khi ngủ do tắc nghẽn chủ yếu do béo phì, đặc điểm hình thái
vùng hàm-mặt và bất thường cấu trúc đường hô hấp trên gây hẹp đường thở khi
ngủ.
% theo nghiên cứu \cite{Young2004_nguyen_nhan}. Tác giả của nghiên cứu cũng đưa
% ra thêm 
Các yếu tố nguy cơ bổ sung gồm hút thuốc, uống rượu, nghẹt mũi, thay đổi nội
tiết tố sau mãn kinh. Yếu tố di truyền cũng góp phần làm tăng khả năng xuất
hiện và mức độ nặng của OSA. Bên cạnh đó, tư thế ngủ được cho rằng đóng vai trò
quan trọng trong việc khởi phát và làm trầm trọng thêm các triệu chứng của OSA.
% \cite{Menon2013_position}. 
Hội chứng ngưng thở khi ngủ do tắc nghẽn dạng phụ thuộc tư thế được mô tả lần
đầu tiên vào năm 1984 ở những người có các đợt ngưng thở và giảm thở xảy ra chủ
yếu khi nằm ngửa.
% \cite{Cartwright1984}.

Để đánh giá chính xác tình trạng OSA, đa ký giấc ngủ
(Polysomnography - PSG) được sử dụng và được xem tiêu chuẩn vàng trong đánh giá giấc ngủ.
Đây là phương pháp ghi đồng thời nhiều tín hiệu sinh lý
gồm điện não đồ (EEG), điện nhãn đồ (EOG), điện cơ đồ (EMG), điện tâm đồ (ECG), nồng độ oxy trong máu cùng các thông số sinh lý liên quan khác.
% \cite{Markun2020_psg}. 
Trong đó, tư thế ngủ luôn là một kênh dữ liệu quan trọng để bác sĩ đưa ra kết
luận về tình trạng OSA và phương pháp điều trị hợp lý. Tuy nhiên, tại Việt Nam,
chi phí cho một lần thực hiện đa ký giấc ngủ còn khá cao (chi phí đi lại, lưu
trú, thăm khám). Bên cạnh đó, việc gắn nhiều điện cực trong khi ngủ gây nhiều
bất tiện cho người bệnh, đồng thời có thể làm sai lệch hoặc gián đoạn tín hiệu
ghi nhận trong quá trình đo.

Trong bối cảnh các công nghệ thiết kế và chế tạo liên tục tiến bộ, kích thước
của vi điều khiển, cảm biến ngày càng được thu nhỏ, đồng thời hiệu năng của
chúng cũng được nâng cao đáng kể. Kết hợp với công nghệ trí tuệ nhân tạo góp
phần nâng cao hiệu quả từ khâu khai thác dữ liệu được thu thập từ cảm biến đến
phân tích dữ liệu đó và có thể đưa ra những hành động mà không cần sự can thiệp
của con người. Đặc biệt, học máy tại biên đang dần phổ biến để giải quyết các
bài toán cần đưa ra quyết định nhanh, nhưng đòi hỏi vẫn giữ độ chính xác cao và
triển khai trên các thiết bị phần cứng có tài nguyên hạn chế đi kèm mức tiêu
thụ năng lượng thấp.

Trên cơ sở đó, được sự đồng ý của Thầy PSG.TS. Mai Anh Tuấn, tôi chọn đề tài
luận văn là “Nghiên cứu, phát triển mô hình học máy tại biên nhằm phân loại tư
thế ngủ”. Luận văn góp phần vào việc giải quyết các vấn đề hết sức cấp bách và
cần thiết trong thực tế.

\noindent\textbf{Mục tiêu của đề tài}

Đề tài hướng tới ba mục tiêu chính bao gồm:
\begin{itemize}
  \item Nghiên cứu, đề xuất hệ thống phần cứng phục vụ đo lường, thu thập và xử lý tín
        hiệu gia tốc, đảm bảo hiệu năng phù hợp cho việc triển khai mô hình học máy tại
        biên trong bài toán phân loại tư thế ngủ ở người.

  \item Nghiên cứu, đề xuất các mô hình học máy thích hợp cho việc phân loại tư thế ngủ
        ở người dựa trên các đặc trưng của tín hiệu gia tốc được trích xuất từ bộ dữ
        liệu cảm biến do nhóm nghiên cứu thu thập.

  \item Chuẩn hóa, thực thi và đánh giá mô hình học máy trên vi điều khiển Nordic
        nRF52840 nhằm kiểm chứng hoạt động thực tế của hệ thống nhận dạng tư thế ngủ
        tại biên.
\end{itemize}

\noindent\textbf{Đối tượng và phạm vi nghiên cứu}

Đối tượng nghiên cứu của đề tài là hội chứng ngưng thở khi ngủ do tắc nghẽn (OSA)
cùn với mối liên hệ với tư thế ngủ, một yếu tố được chứng minh có khả năng làm
khởi phát hoặc làm trầm trọng thêm mức độ tắc nghẽn đường thở.
Bên cạnh việc xem xét đặc điểm lâm sàng của OSA, các kênh tín hiệu đa kí giấc ngủ,
nghiên cứu tập trung
vào tín hiệu thu nhận từ cảm biến
gia tốc gắn tại vùng cổ nơi cho thấy độ ổn định
cao trong giám sát tư thế ngủ.
Sau đó, tiến tới nghiên cứu về các đặc trưng, các mô hình phù hợp với bộ dữ liệu đó để phân loại tư thế ngủ.

Phạm vi nghiên cứu: Dữ liệu thu thập từ 25 tình nguyện viên dưới sự giám sát
của nhóm và camera. Các mô hình học máy truyển thống: rừng ngẫu nhiên (Random
Forest - RF), máy véc-tơ hỗ trợ (Support Vector Machine - SVM), hồi quy
Logistic (Logistic Regression - LR), Gradient Boosting (GB) và mô hình học sâu
(Artificial Neural Network - ANN). Triển khai trên phần cứng được thiết kế từ
vi điều khiển nrf52840.

\noindent\textbf{Phương pháp nghiên cứu}

Nghiên cứu này được thực hiện dựa trên kết hợp hai phương pháp giữa nghiên cứu
lý thuyết bằng cách tổng hợp tài liệu và nghiên cứu thực nghiệm.

Đầu tiên, tổng
hợp lý thuyết các vấn đề liên quan trong luận văn từ giáo trình, bài giảng,
luận văn, các công trình khoa học công bố có liên quan. Sau đó, phân tích, cô
đọng để viết một bài báo tổng quan về công nghệ chẩn đoán OSA tại nhà.

Thứ hai, tiến hành phát triển hệ thống thu thập, lưu trữ (phần cứng, phần mềm)
dựa trên nguyên tắc điện toán đám mây. Sau đó, đánh giá đặc trưng, chọn lọc
tham số mô hình dựa trên lý thuyết đã tìm hiểu. Kết quả sẽ được đánh giá, so
sánh với các công trình có liên quan.

\noindent\textbf{Bố cục luận văn}

Cấu trúc luận văn được trình bày trong bốn chương chính như sau:

\noindent\textbf{Chương 1:} Tổng quan OSA, mối liên hệ với tư thế ngủ và kĩ thuật sàng
lọc, phân loại.

\noindent\textbf{Chương 2:} Hệ thống thu thập, xử lý tín hiệu cảm biến và đánh giá bằng
mô hình học máy.

\noindent\textbf{Chương 3:} Kết quả và đánh giá.

% \changefontsizes[16pt]{13pt}

\vspace{0.2cm}
\noindent\textbf{\textit{Lưu ý:}} Trong toàn bộ luận văn, các giá trị thập phân được biểu diễn bằng dấu “.” và các giá trị hàng nghìn được phân tách bằng dấu “,”.
